\documentclass[a4paper, french, 11 pt]{article}

% Language setting
% Replace `english' with e.g. `french' to change the document language
\usepackage[french]{babel}

% Set page size and margins
% Replace `letterpaper' with `a4paper' for UK/EU standard size
\usepackage[letterpaper,top=2cm,bottom=2cm,left=2cm,right=2cm,marginparwidth=1.75cm]{geometry}

% Useful packages
\usepackage{graphicx}
\usepackage{siunitx}
\usepackage[colorlinks=true, allcolors=blue]{hyperref}
%\usepackage{pgfplots}
%\usepackage{tkz-euclide}
\usepackage[T1]{fontenc}
\usepackage[utf8]{inputenc}
%\usepackage[european, straightvoltages, RPvoltages]{circuitikz}
%\usetikzlibrary{babel}
\usepackage{indentfirst}
%\usepackage{esint}
\usepackage{fancyhdr}
%\\
\usepackage{amsfonts, amsmath, amssymb}
\usepackage{subcaption}
%\usepackage[export]{adjustbox}
%\usepackage{lipsum}
%\usepackage{mathtools}
%\usepackage{tkz-tab}
%\usepackage{longtable}
%\usepackage[cyr]{aeguill}
%\usepackage{xspace}
%\usepackage{titling}
%\usepackage[absolute,overlay]{textpos}
%\usepackage{etoc}
\usepackage{enumitem}
\usepackage{multirow}
\usepackage{nicefrac}
\usepackage{mathabx}
\usepackage[autostyle=true, french=guillemets]{csquotes}
%\pgfplotsset{compat=newest}

%Mathematical operators
\DeclareMathOperator{\e}{e}
\newcommand{\deriv}{\mathrm{d}}

%Bibliography
\usepackage[
backend=biber,
style=numeric-comp,
sorting=nty
]{biblatex}
\addbibresource{Reference_stage.bib}

%Style
\pagestyle{fancy}
\fancyhf{}
\fancyhead[C]{\textsc{ }}
\fancyhead[R]{\textsc{Stage Magistère}}
\fancyfoot[C]{\thepage}

\title{\textbf{Stage Magistère LPSC}\\ Méthodologie de recherche}
\author{Tibor \textsc{Donin de Rosière}}



\begin{document}
\maketitle
\section{Formulation de la problématique}
\begin{itemize}[label=\textbullet]
\item Implémentation d'un modèle de gaz chaud en partant des modèles de halos de matière noire
\item Étendre le cadre de distribution de l'occupation du halo à la prévision du gaz chaud
\item Mettre en œuvre un modèle de gaz chaud dans le cadre de la distribution d'occupation du halo
\item Élaborer un modèle gén,éral décrivant la répartition des composants de la matière concernée ainsi que de la thermodynamique des gaz dans le cadre du modèle des halos de structure à grande échelle.
\end{itemize}

\section{Analyse de la problématique et questions}
Quel est l'impact des trous noirs et des étoiles sur la répartition du gaz au sein d'un halo ?\\

\section{Protocole}
\subsection{Système simplifié}
Distribution de matière (gaz, étoile, matière noire) sous forme de halo, dont les échelles typiques de taille sont de l'ordre du Mpc. Au centre du halo, il y a une galaxie (10 pc à 30 kpc) composée essentiellement de matière baryonique (étoiles, nuages de gaz, poussières). Le noyau galactique actif et les supernovae libèrent de grandes quantités de gaz sous forme de vent galactique. Cela expulse du gaz chaud en dehors de l'halo. Aux extrémités de l'halo, il y a une répartition de gaz plus diffus à l'équilibre hydrostatique, qui est ensuite pertubé par les jets de gaz chaud \cite{pandey2024godmaxmodelinggasthermodynamics}.\\
Les galaxies satellites sont décentrés par rapport à leurs halos hôtes et se trouvent généralement dans des hôtes plus massifs qu'un hôte dont le centre présente une masse stellaire similaire \cite{mccarthy2025flamingocombiningkineticsz}.\\
Les flux sortants transportent le gaz du centre du halo vers la périphérie, ce qui rend le gaz moins dense au centre et augmente l'intensité de la rétroaction. On peut supposer que, en raison de la structure anisotrope des flux sortants, la redistribution du gaz est elle aussi anisotrope, transportant davantage de gaz vers les régions de faible densité que vers celles de forte densité, et rendant la distribution globale du gaz plus sphérique \cite{Ondaro_Mallea_2025}.\\
Schéma

\newpage
\subsection{Grandeurs pertinentes}

\begin{table}[h]
\centering
\begin{tabular}{|l c c|}
    \hline
    Paramètre & Symbole & Intervalle de valeurs \\
    \hline
    Densité massique du bulbe galactique & $\rho_{cga}$ & \\
    Densité des composantes du vent stellaire & & \\
    Densité massique du gaz & $\rho_g$ & $[10^{11}\ ;\ 10^{13}]\ \unit{M_{\Sun}.Mpc^{-3}}$ \cite{schneider2025baryonificationalternativehydrodynamicalsimulations} \\
    Masse du bulbe galactique & & \\
    Masse du gaz & $m_g$ & $[10^{8}\ ;\ 10^{9}]\ \unit{M_{\Sun}}$\footnote{Valeurs d'essai sur un volume de (1 Gpc)$^3$} \cite{kovac2025baryonificationiiconstrainingfeedback}\\
    Masse du halo & $M_h$ & $[10^{12}\ ;\ 10^{15}]$ M$_{\Sun}$\\
    Pression du vent stellaire & & \\
    Pression du gaz & & \\
    Redshift & $z$ & $0-3$ \\
    Spectre de puissance de la matière linéaire & $P$ & \\
    Température du vent stellaire & & \\
    Température du gaz du halo galactique & $T_g$ & $[10^{6}\ ;\ 10^8]\ \unit{K}$\footnote{Calculée à partir de l'équation \ref{eq:T200} pour $z=0$}\\
    Vitesse moyenne des flux sortants & $v_{jet}$ & $[1\ ;\ 65]\ \unit{km.s^{-1}}$ \cite{Ondaro_Mallea_2025}\\
    \hline
\end{tabular} 
\caption{}
\end{table}
Pour du gaz autour de halos isolés de trois classes de masse à $z = 0$, on a 
$\rho_g\in [10^{-1}\ ;\ 10^{3}]\ \unit{M_{\Sun}.Mpc^{-3}}$ \cite{Ondaro_Mallea_2025}
\subsection{Lois fondamentales}
Les disciplines mises en jeux sont de la thermodynamique et de l'hydrodynamique.\\
Equation de l'équilibre hydrostatique :
\begin{equation}
\frac{\deriv P_{tot}}{\deriv r}=-\rho_g(r)\frac{GM_{dmb}(r)}{r^2}
\end{equation} 
Profil de densité de matière noire : NFW tronquée\\
\begin{equation}
\rho_{\text{nfw}}=\frac{\rho_{\text{nfw,0}}}{x_s(1+x_s)^2}\frac{1}{(1+x_t^2)^2}
\end{equation}
Fraction baryonique :
\begin{equation}
f_\text{b}=\frac{\Omega_\text{b}}{\Omega_\text{m}}
\end{equation}
L'effet Sunyaev-Zeldovich est pris en compte dans le calcul des composantes.\\

L'effet thermique de Sunyaev-Zeldovich\\
Facteur de Compton
\begin{equation}
y = \int \sigma_T\, n_e\, \frac{k_B T_e}{m_e c^2}\, \deriv l
= \frac{\sigma_T}{m_e c^2}\int P_e\, \deriv l
\end{equation}
L'effet cinématique de Sunyaev-Zeldovich est une anisotropie secondaire du fond diffus cosmologique (CMB). Il est dû à la diffusion Compton inverse des photons du CMB avec les électrons libres présents dans les galaxies et les amas qui présentent un mouvement global non nul. Cette diffusion inverse modifie la température thermodynamique du CMB selon la formule suivante \cite{kovac2025baryonificationiiconstrainingfeedback} \cite{siegel2025jointxraykineticsunyaevzeldovich}
\begin{equation}
\frac{\Delta T_{\mathrm{kSZ}}}{T_{\mathrm{CMB}}}=\tau\left(\frac{v_\mathrm{r}}{c}\right).
\end{equation}
où $\tau$ est la profondeur optique due à la diffusion de Thompson entre l'observateur et le redshift $z$. Pour plus de précisions sur cette effet, voir \cite{bigwood2025kineticsunyaevzeldovicheffect}\\
Le spectre de puissance kSZ est une intégrale pondérée de $\Delta_B^2$ par rapport au décalage vers le rouge
Le profil de densité des électrons libres peut être déterminé comme suit :
\begin{equation}
n_e(r)=\frac{\rho_\text{g}(r)}{m_p\mu_e}
\end{equation}
On en déduit l'expression de la pression des électrons libres :
\begin{equation}
P_e=k_\text{B}n_eT_e
\end{equation}
Expression de la température de l'halo $T_{200}$ à partir de $M_{200}$ et $R_{200}$ \cite{Battaglia_2012}
\begin{equation}
\label{eq:T200}
kT_{200} = \frac{G M_{200} \ \mu\ m_\text{p}}{3 R_{200}}
\end{equation}
Luminosité de surface des rayons X :
\begin{equation}
S_X\propto\int n_e^2\Lambda(T_e)\deriv l
\end{equation}
où $\Lambda(T_e)$ est la fonction de refroidissement par rayons X. Une modélisation du coefficient d'émissivité volumique dans le domaine des rayons X faibles (0,5–2 keV) \cite{Oppenheimer_2025} :
\begin{equation}
\Lambda(T, Z) = \Lambda_0
    \left(\frac{T}{1\,\text{keV}}\right)^{\alpha_T}
    \left(\frac{Z}{0.3\,Z_\odot}\right)^{\alpha_Z}
\end{equation}
Le spectre de puissance linéaire de la matière est calculée avec CAMB \cite{Lewis_2000} 
%\Delta^2(k, z) \equiv \frac{k^3 P(k, z)}{2\pi^2}
\begin{equation}
P_{\rm lin}(k, z)\ \propto\ k^{n_s} T^2(k)\, D^2(z)
\quad [(h^{-1}\,{\rm Mpc})^3]
\end{equation}
où $T(k)$ est la fonction de transmission de matière (matter transfer function) et $D(z)$ le facteur d'accroissement linéaire (linear growth factor) définit tel que
\begin{equation}
D(z) = \frac{5\Omega_m}{2} H(z) / H_0
\int_z^\infty \frac{(1+z')}{[H(z')/H_0]^3}\,\deriv z'.
\end{equation}
\subsection{Modélisation}
Dans les travaux du dépôt git \verb|hod_mod|, il n'y a aucun modèle et aucune simulation de l'effet cinématique Sunyaev-Zelodvich, et celui-ci traduit l'effet des électrons libres présent dans les halos sur le fond diffus cosmologique. Cette diffusion inverse influe la température du CMB comme
\begin{equation}
    \frac{\Delta T_\mathrm{kSZ}}{T_\mathrm{CMB}} = \frac{\sigma_\mathrm{T}}{c} \int_\mathrm{los} \mathrm{d}l\,e^{-\tau} n_\mathrm{e} \; v_\mathrm{p}
\end{equation}
Nous supposons que le gaz se déplace de manière cohérente avec la vitesse globale du halo (voir réf. \cite{Hadzhiyska2023}), ce qui permet d'éliminer la vitesse particulière $v_p$ de l'intégrale dans la ligne de visée. Cette simplification est couramment utilisée dans les analyses par empilement, où les vitesses individuelles ne sont pas connues avec précision. Au lieu de cela, le signal attendu est modélisé à l'aide de la vitesse radiale quadratique moyenne (RMS) de l'échantillon de halos, que nous fixons à $v_r = 1,06 \times 10^{-3}c$, conformément à la référence \cite{Amodeo2021}. La profondeur optique $\tau$ le long de la ligne de visée est définit comme suit :
\begin{equation}
\tau(\theta) =  \sigma_\mathrm{T} \int_{\rm los} n_{\mathrm{e}}\left(\sqrt{l^2\:+\:d_A(z)^2\theta^2}\right)\deriv l\;,
\end{equation}

\section{Résultats notables existants}
À grande échelle, la vitesse du gaz est identique à celle de la
matière noire. Cependant, le champ de densité du gaz est sous-estimé par rapport à celui de la matière noire, même à des échelles aussi grandes que $k = 0,01\ \unit{h/Mpc}$. (à chercher le sens physique)\\
Le biais est pratiquement indépendant de la fraction de gaz dans les amas, mais il dépend sensiblement de la fonction de masse stellaire. Ceci est interprété comme une conséquence de la formation d'étoiles dans les régions denses et agglomérées, qui laisse le gaz se concentrer de préférence dans les régions moins denses.\\
À petites échelles, le gaz est fortement influencé par des forces non gravitationnelles.
Le spectre de puissance du gaz est nettement atténué par rapport au cas où seule la gravité est prise en compte, tout comme les écoulements potentiels et les vitesses moyennes par paires. Bien qu'il s'agisse sans doute de l'effet combiné de plusieurs processus non gravitationnels, ce phénomène est étroitement lié à l'intensité de la rétroaction dans les simulations FLAMINGO.\\
La rétroaction des noyaux galactiques actifs (AGN) remodèle le champ de densité en repoussant le gaz loin des halos et elle est particulièrement efficace dans les halos à l'échelle des groupes dont la masse $M_{200}$ est d'environ $10^{13}\ \unit{h^{-1} M_{\Sun}}$\\
Les variations des profils de densité et de vitesse radiale induites par le retour d'énergie peuvent s'expliquer principalement par d'importants écoulements provenant de l'accrétion sur des trous noirs supermassifs. Que le retour d'énergie de l'AGN soit modélisé sous forme thermique ou de jets, le gaz environnant finit par se dilater, créant des bulles flottantes pouvant atteindre des dizaines de rayons viriaux avec des vitesses radiales comparables à la vitesse d'échappement du halo. \cite{Ondaro_Mallea_2025}\\
\medskip
\printbibliography[
title={Références}
]
\end{document}